\IfFileExists{stacks-project.cls}{%
\documentclass{stacks-project}
}{%
\documentclass{amsart}
}
\usepackage{amssymb}

% For dealing with references we use the comment environment
\usepackage{verbatim}
\usepackage{fancyvrb}
\usepackage{pmboxdraw}
\newenvironment{reference}{\comment}{\endcomment}
%\newenvironment{reference}{}{}
\newenvironment{slogan}{\comment}{\endcomment}
\newenvironment{history}{\comment}{\endcomment}

% For source figures retained as assets
\usepackage{graphicx}

% For commutative diagrams we use Xy-pic
\usepackage[all]{xy}

% We use 2cell for 2-commutative diagrams.
\xyoption{2cell}
\UseAllTwocells

% We use multicol for the list of chapters between chapters
\usepackage{multicol}

% This is generally recommended for better output
\usepackage{lmodern}
\usepackage[T1]{fontenc}
\usepackage{textalpha}

% For cross-file-references
% We prefer the snapshot tmp/xref/#2.aux,
% which latexmkrc replaces atomically, so
% that documents compiled simultaneously
% never read a half-written .aux file.
% The last argument keeps links to #2.pdf.
\usepackage{xr-hyper}
\newcommand{\maybeexternaldocument}[2]{%
  \IfFileExists{#2.aux}{%
    \externaldocument[#1]{#2}%
  }{%
    \PackageWarning{stack}{Missing #2.aux; external references with prefix #1 unavailable}%
  }%
}
\let\externaldocumentorig\externaldocument
\renewcommand{\externaldocument}[2][]{%
  \IfFileExists{tmp/xref/#2.aux}{%
    \externaldocumentorig[#1]{tmp/xref/#2}[#2.pdf]%
  }{%
  \IfFileExists{#2.aux}{%
    \externaldocumentorig[#1]{#2}%
  }{%
    \PackageWarning{stack}{Missing #2.aux; external references with prefix #1 unavailable}%
  }%
  }%
}

% Package for hypertext links:
\usepackage{hyperref}

% For any local file, say "hello.tex" you want to link to please
% use \externaldocument[hello-]{hello}
\externaldocument[introduction-]{introduction}
\externaldocument[conventions-]{conventions}
\externaldocument[sets-]{sets}
\externaldocument[categories-]{categories}
\externaldocument[topology-]{topology}
\externaldocument[sheaves-]{sheaves}
\externaldocument[sites-]{sites}
\externaldocument[stacks-]{stacks}
\externaldocument[fields-]{fields}
\externaldocument[algebra-]{algebra}
\externaldocument[brauer-]{brauer}
\externaldocument[homology-]{homology}
\externaldocument[derived-]{derived}
\externaldocument[simplicial-]{simplicial}
\externaldocument[more-algebra-]{more-algebra}
\externaldocument[smoothing-]{smoothing}
\externaldocument[modules-]{modules}
\externaldocument[sites-modules-]{sites-modules}
\externaldocument[injectives-]{injectives}
\externaldocument[cohomology-]{cohomology}
\externaldocument[sites-cohomology-]{sites-cohomology}
\externaldocument[dga-]{dga}
\externaldocument[dpa-]{dpa}
\externaldocument[sdga-]{sdga}
\externaldocument[hypercovering-]{hypercovering}
\externaldocument[schemes-]{schemes}
\externaldocument[constructions-]{constructions}
\externaldocument[properties-]{properties}
\externaldocument[morphisms-]{morphisms}
\externaldocument[coherent-]{coherent}
\externaldocument[divisors-]{divisors}
\externaldocument[limits-]{limits}
\externaldocument[varieties-]{varieties}
\externaldocument[topologies-]{topologies}
\externaldocument[descent-]{descent}
\externaldocument[perfect-]{perfect}
\externaldocument[more-morphisms-]{more-morphisms}
\externaldocument[flat-]{flat}
\externaldocument[groupoids-]{groupoids}
\externaldocument[more-groupoids-]{more-groupoids}
\externaldocument[etale-]{etale}
\externaldocument[chow-]{chow}
\externaldocument[intersection-]{intersection}
\externaldocument[pic-]{pic}
\externaldocument[weil-]{weil}
\externaldocument[adequate-]{adequate}
\externaldocument[dualizing-]{dualizing}
\externaldocument[duality-]{duality}
\externaldocument[discriminant-]{discriminant}
\externaldocument[derham-]{derham}
\externaldocument[local-cohomology-]{local-cohomology}
\externaldocument[algebraization-]{algebraization}
\externaldocument[curves-]{curves}
\externaldocument[resolve-]{resolve}
\externaldocument[models-]{models}
\externaldocument[functors-]{functors}
\externaldocument[equiv-]{equiv}
\externaldocument[pione-]{pione}
\externaldocument[etale-cohomology-]{etale-cohomology}
\externaldocument[proetale-]{proetale}
\externaldocument[relative-cycles-]{relative-cycles}
\externaldocument[more-etale-]{more-etale}
\externaldocument[trace-]{trace}
\externaldocument[crystalline-]{crystalline}
\externaldocument[spaces-]{spaces}
\externaldocument[spaces-properties-]{spaces-properties}
\externaldocument[spaces-morphisms-]{spaces-morphisms}
\externaldocument[decent-spaces-]{decent-spaces}
\externaldocument[spaces-cohomology-]{spaces-cohomology}
\externaldocument[spaces-limits-]{spaces-limits}
\externaldocument[spaces-divisors-]{spaces-divisors}
\externaldocument[spaces-over-fields-]{spaces-over-fields}
\externaldocument[spaces-topologies-]{spaces-topologies}
\externaldocument[spaces-descent-]{spaces-descent}
\externaldocument[spaces-perfect-]{spaces-perfect}
\externaldocument[spaces-more-morphisms-]{spaces-more-morphisms}
\externaldocument[spaces-flat-]{spaces-flat}
\externaldocument[spaces-groupoids-]{spaces-groupoids}
\externaldocument[spaces-more-groupoids-]{spaces-more-groupoids}
\externaldocument[bootstrap-]{bootstrap}
\externaldocument[spaces-pushouts-]{spaces-pushouts}
\externaldocument[spaces-chow-]{spaces-chow}
\externaldocument[groupoids-quotients-]{groupoids-quotients}
\externaldocument[spaces-more-cohomology-]{spaces-more-cohomology}
\externaldocument[spaces-simplicial-]{spaces-simplicial}
\externaldocument[spaces-duality-]{spaces-duality}
\externaldocument[formal-spaces-]{formal-spaces}
\externaldocument[restricted-]{restricted}
\externaldocument[spaces-resolve-]{spaces-resolve}
\externaldocument[formal-defos-]{formal-defos}
\externaldocument[defos-]{defos}
\externaldocument[cotangent-]{cotangent}
\externaldocument[examples-defos-]{examples-defos}
\externaldocument[algebraic-]{algebraic}
\externaldocument[examples-stacks-]{examples-stacks}
\externaldocument[stacks-sheaves-]{stacks-sheaves}
\externaldocument[criteria-]{criteria}
\externaldocument[artin-]{artin}
\externaldocument[quot-]{quot}
\externaldocument[stacks-properties-]{stacks-properties}
\externaldocument[stacks-morphisms-]{stacks-morphisms}
\externaldocument[stacks-limits-]{stacks-limits}
\externaldocument[stacks-cohomology-]{stacks-cohomology}
\externaldocument[stacks-perfect-]{stacks-perfect}
\externaldocument[stacks-introduction-]{stacks-introduction}
\externaldocument[stacks-more-morphisms-]{stacks-more-morphisms}
\externaldocument[stacks-geometry-]{stacks-geometry}
\externaldocument[moduli-]{moduli}
\externaldocument[moduli-curves-]{moduli-curves}
\externaldocument[examples-]{examples}
\externaldocument[exercises-]{exercises}
\externaldocument[guide-]{guide}
\externaldocument[desirables-]{desirables}
\externaldocument[coding-]{coding}
\externaldocument[obsolete-]{obsolete}
\externaldocument[fdl-]{fdl}
\externaldocument[index-]{index}

% Heap Project documents
\externaldocument[k3-]{k3}
\externaldocument[torelli-]{torelli}
\externaldocument[linear-algebra-]{linear-algebra}
\externaldocument[groups-]{groups}
\externaldocument[complex-analysis-]{complex-analysis}
\externaldocument[functional-analysis-]{functional-analysis}
\externaldocument[categories-]{categories}
\externaldocument[commutative-algebra-]{commutative-algebra}
\externaldocument[ringed-spaces-]{ringed-spaces}
\externaldocument[homological-algebra-]{homological-algebra}
\externaldocument[lie-algebras-]{lie-algebras}
\externaldocument[representation-theory-]{representation-theory}
\externaldocument[smooth-manifolds-]{smooth-manifolds}
\externaldocument[differential-topology-]{differential-topology}
\externaldocument[algebraic-topology-]{algebraic-topology}
\externaldocument[differential-geometry-]{differential-geometry}
\externaldocument[algebraic-geometry-]{algebraic-geometry}
\externaldocument[symplectic-geometry-]{symplectic-geometry}
\externaldocument[complex-geometry-]{complex-geometry}
\externaldocument[hodge-theory-]{hodge-theory}
\externaldocument[%
variations-of-hodge-structures-]{%
variations-of-hodge-structures}
\externaldocument[%
homological-mirror-symmetry-]{%
homological-mirror-symmetry}
\maybeexternaldocument{deformations-}{deformations}
\externaldocument[vertex-algebras-]{vertex-algebras}
\externaldocument[springer-]{springer}
\externaldocument[infty-categories-]{infty-categories}
\externaldocument[seminars-differential-geometry-]{%
seminars-differential-geometry}
\externaldocument[seminars-complex-geometry-]{%
seminars-complex-geometry}
\externaldocument[seminars-algebraic-geometry-]{%
seminars-algebraic-geometry}
\externaldocument[seminars-analysis-pde-]{%
seminars-analysis-pde}
\externaldocument[surfaces-]{surfaces}
\externaldocument[birational-maps-commutative-algebra-]{birational-maps-commutative-algebra}
\externaldocument[cremona-transformations-]{cremona-transformations}
\externaldocument[derived-categories-of-sheaves-]{derived-categories-of-sheaves}
\externaldocument[geometric-invariant-theory-]{geometric-invariant-theory}
\externaldocument[geometric-stability-conditions-and-group-actions-]{geometric-stability-conditions-and-group-actions}
\externaldocument[moduli-spaces-of-sheaves-]{moduli-spaces-of-sheaves}
\externaldocument[bridgeland-stability-general-theory-]{bridgeland-stability-general-theory}
\externaldocument[hilbert-schemes-]{hilbert-schemes}
\externaldocument[noncommutative-abelian-surfaces-and-kummer-type-hyperkahler-varieties-]{noncommutative-abelian-surfaces-and-kummer-type-hyperkahler-varieties}
\externaldocument[mumford-tate-groups-in-hodge-theory-]{mumford-tate-groups-in-hodge-theory}
\externaldocument[hodge-birational-atoms-2026-]{hodge-birational-atoms-2026}
\externaldocument[xxii-egd-2026-]{%
xxii-egd-2026}
\externaldocument[pde-]{pde}
\externaldocument[probability-]{probability}

\externaldocument[linguistics-]{linguistics}

\externaldocument[%
16th-alga-meeting-2026-]{%
16th-alga-meeting-2026}

\externaldocument[%
iv-jornadas-lefschetz-2026-]{%
iv-jornadas-lefschetz-2026}

\externaldocument[reading-the-masters-]{%
reading-the-masters}

\externaldocument[real-analysis-]{%
real-analysis}

\externaldocument[optimal-transport-]{%
optimal-transport}

\externaldocument[history-]{history}

\externaldocument[seiberg-witten-]{seiberg-witten}

\externaldocument[history-philosophy-biology-]{history-philosophy-biology}

\externaldocument[classical-mechanics-]{classical-mechanics}
\externaldocument[electromagnetism-]{electromagnetism}
\externaldocument[waves-and-optics-]{waves-and-optics}
\externaldocument[wif-]{waves-in-fluids-iv}
\externaldocument[thermal-physics-]{thermal-physics}
\externaldocument[quantum-mechanics-]{quantum-mechanics}
\externaldocument[special-relativity-]{special-relativity}
\externaldocument[physics-]{physics}

\externaldocument[chemistry-]{chemistry}

\externaldocument[molecular-cell-biology-]{molecular-cell-biology}

\externaldocument[microbiology-]{microbiology}

\externaldocument[protistology-]{protistology}

\externaldocument[mycology-]{mycology}

\externaldocument[botany-]{botany}

\externaldocument[zoology-]{zoology}

\externaldocument[animal-histology-]{animal-histology}

\externaldocument[animal-physiology-]{animal-physiology}

\externaldocument[systematics-biogeography-]{systematics-biogeography}

\externaldocument[ecology-natural-resources-]{ecology-natural-resources}

\externaldocument[earth-sciences-]{earth-sciences}

\externaldocument[biostatistics-]{biostatistics}

\externaldocument[mathematical-systems-biology-]{mathematical-systems-biology}

\externaldocument[cimpa-el-salvador-2024-]{cimpa-el-salvador-2024}

\externaldocument[anthropology-]{anthropology}

\externaldocument[yang-mills-theory-]{yang-mills-theory}

\externaldocument[seminars-algebra-]{seminars-algebra}

\externaldocument[seminars-physics-]{seminars-physics}

\externaldocument[seminars-computing-]{seminars-computing}

\externaldocument[oktobermat-2026-]{oktobermat-2026}

\externaldocument[saet-2026-]{saet-2026}

\externaldocument[seminars-symplectic-geometry-]{seminars-symplectic-geometry}

\externaldocument[seminars-dynamics-]{seminars-dynamics}

\externaldocument[oktobermat-2024-]{oktobermat-2024}

\externaldocument[quantum-groups-]{quantum-groups}

\externaldocument[legal-2027-]{legal-2027}

\externaldocument[derived-categories-]{derived-categories}

% Theorem environments.
%
\theoremstyle{plain}
\newtheorem{theorem}[subsection]{Theorem}
\newtheorem{proposition}[subsection]{Proposition}
\newtheorem{lemma}[subsection]{Lemma}
\newtheorem{corollary}[subsection]{Corollary}

\theoremstyle{definition}
\newtheorem{definition}[subsection]{Definition}
\newtheorem{example}[subsection]{Example}
\newtheorem{exercise}[subsection]{Exercise}
\newtheorem{situation}[subsection]{Situation}

\theoremstyle{remark}
\newtheorem{remark}[subsection]{Remark}
\newtheorem{remarks}[subsection]{Remarks}

\numberwithin{equation}{subsection}

% Macros
%
\def\lim{\mathop{\mathrm{lim}}\nolimits}
\def\colim{\mathop{\mathrm{colim}}\nolimits}
\def\Spec{\mathop{\mathrm{Spec}}}
\def\Hom{\mathop{\mathrm{Hom}}\nolimits}
\def\Ext{\mathop{\mathrm{Ext}}\nolimits}
\def\SheafHom{\mathop{\mathcal{H}\!\mathit{om}}\nolimits}
\def\SheafExt{\mathop{\mathcal{E}\!\mathit{xt}}\nolimits}
\def\Sch{\mathit{Sch}}
\def\Mor{\mathop{\mathrm{Mor}}\nolimits}
\def\Ob{\mathop{\mathrm{Ob}}\nolimits}
\def\Sh{\mathop{\mathit{Sh}}\nolimits}
\def\NL{\mathop{N\!L}\nolimits}
\def\CH{\mathop{\mathrm{CH}}\nolimits}
\def\proetale{{pro\text{-}\acute{e}tale}}
\def\etale{{\acute{e}tale}}
\def\QCoh{\mathit{QCoh}}
\def\Ker{\mathop{\mathrm{Ker}}}
\def\Im{\mathop{\mathrm{Im}}}
\def\Coker{\mathop{\mathrm{Coker}}}
\def\Coim{\mathop{\mathrm{Coim}}}

% Boxtimes
%
\DeclareMathSymbol{\boxtimes}{\mathbin}{AMSa}{"02}

%
% Macros for moduli stacks/spaces
%
\def\QCohstack{\mathcal{QC}\!\mathit{oh}}
\def\Cohstack{\mathcal{C}\!\mathit{oh}}
\def\Spacesstack{\mathcal{S}\!\mathit{paces}}
\def\Quotfunctor{\mathrm{Quot}}
\def\Hilbfunctor{\mathrm{Hilb}}
\def\Curvesstack{\mathcal{C}\!\mathit{urves}}
\def\Polarizedstack{\mathcal{P}\!\mathit{olarized}}
\def\Complexesstack{\mathcal{C}\!\mathit{omplexes}}
% \Pic is the operator that assigns to X its picard group, usage \Pic(X)
% \Picardstack_{X/B} denotes the Picard stack of X over B
% \Picardfunctor_{X/B} denotes the Picard functor of X over B
\def\Pic{\mathop{\mathrm{Pic}}\nolimits}
\def\Picardstack{\mathcal{P}\!\mathit{ic}}
\def\Picardfunctor{\mathrm{Pic}}
\def\Deformationcategory{\mathcal{D}\!\mathit{ef}}


\begin{document}

\title{Birman-Hilden theory}
\maketitle

\phantomsection
\label{section-phantom}

\tableofcontents

\section{Overview}
\label{section-introduction}

\noindent
These are my notes from a seminar on
Birman-Hilden theory. They were written
in Spanish under the title
{\it Teor\'ia de Birman-Hilden}; the
original is kept in {\tt standalone/} as
{\tt teoria-birman-hilden}. The source
does not name the speakers or the dates
of the sessions. This chapter is a
translation that keeps the content as it
was: statements without proof stay
without proof, unfinished proofs stay
unfinished, and questions stay open.
Where the original is wrong, incomplete
or unclear, a sentence starting with
{\it Note on the original} says so.
Statements marked ``Homework'' were
marked {\it Tarea} in the original.

\medskip\noindent
{\bf Big picture.} A branched cover
$p:S\to X$ of surfaces relates
homeomorphisms of $S$ that preserve the
fibers of $p$ to homeomorphisms of $X$
that lift. Birman-Hilden theory asks
when this relation passes to isotopy
classes. The notes go through the
following steps.
\begin{enumerate}
\item Covering spaces, group actions,
lifts and deck transformations
(Sections
\ref{section-covering-spaces},
\ref{section-group-actions},
\ref{section-lifts} and
\ref{section-deck-transformations}).
\item Covering actions and the
quotients they produce
(Section \ref{section-covering-actions}).
\item Branched covers and the
Riemann-Hurwitz formula
(Section \ref{section-branched-covers}).
\item The mapping class group and first
examples
(Section
\ref{section-mapping-class-group}).
\item The Birman-Hilden property, the
statement of the theorem of Birman and
Hilden, and lemmas towards a proof
(Section
\ref{section-birman-hilden-theorem}).
\end{enumerate}

\section{Covering spaces}
\label{section-covering-spaces}

\noindent
In the proof of the fundamental group of
the circle we had already used the idea
of a covering space. Here is the general
definition.

\begin{definition}
\label{definition-covering-space}
A {\it covering space} of $X$ is a space
$E$ together with a map
$$
p:E\to X
$$
such that every point $x\in X$ has an
open neighbourhood $U\subseteq X$ such
that $p^{-1}(U)$ is a disjoint union of
open sets of $E$, each of which is
projected homeomorphically onto $U$ by
$p$.
\end{definition}

The original has a figure here: a
rectangle $X$ containing a disc $U$, and
above it a rectangle $E$ containing a
stack of five discs, with an arrow $p$
from $E$ down to $X$.

We introduce other terms.
\begin{itemize}
\item The map $p$ is called the
{\it covering projection}.
\item The set $U$ is called a
{\it regular neighbourhood}.
\item The open sets of $p^{-1}(U)$ that
are projected homeomorphically onto $U$
are called the {\it sheets} of $E$ over
$U$.
\item The set $p^{-1}(x)$ is called the
{\it fiber} of $x$ under $p$.
\end{itemize}

\begin{remark}
\label{remark-sheets}
\begin{itemize}
\item If $U$ is connected then the
sheets of $U$ are the connected
components of $p^{-1}(U)$.
\item The number of sheets over $U$ is
equal to the cardinality of $p^{-1}(x)$.
\item This number is {\it locally
constant}, that is, it is the same for
every point.
\end{itemize}
Note on the original: the last item
ends with a colon and nothing follows
it. Locally constant means constant on
a neighbourhood of each point; the
number is the same for every point when
$X$ is connected.
\end{remark}

\begin{example}
\label{example-covering-spaces}
\begin{enumerate}
\item The identity
$\operatorname{Id}:X\to X$.
\item The natural projection of the
disjoint union of copies of a space,
that is,
$p:\bigsqcup_{\alpha\in I}X\to X$.
\item For any two covers $p:E\to X$ and
$q:\tilde{E}\to X$, the projection from
the disjoint union
$r:E\sqcup\tilde{E}\to X$ defined in the
natural way is also a cover. This is to
say that we are interested in covers of
connected spaces whose covering spaces
are connected.
\item The map $p:\mathbb{R}\to S^1$ as
in the proof of the fundamental group of
the circle.
\item For a natural number $n$, the map
$p:S^1\to S^1$ given by $p(z)=z^n$.
\end{enumerate}
\end{example}

\section{Actions of a group on a
topological space}
\label{section-group-actions}

\begin{definition}
\label{definition-action}
We say that a group $G$ {\it acts} on
$X$ if there is a group homomorphism
$$
\varphi:G\to\operatorname{Homeo}(X)
$$
that sends $g$ to the homeomorphism
$\varphi(g):X\to X$, which in turn sends
$x$ to $\varphi(g)(x)$, so that
\begin{itemize}
\item $\varphi(g_1g_2)
=\varphi(g_1)\varphi(g_2)$,
\item $(g_1g_2)x=g_1(g_2(x))$,
\item $\varphi(1_G)
=\operatorname{Id}_X$.
\end{itemize}
This is denoted by
$G\curvearrowright X$. An action is
{\it free} if $g(y)=y$, for any $g\in G$
and $y\in X$, implies that $g=1_G$; and
an action is {\it transitive} if for any
$y,y'\in X$ there is an element $g\in G$
such that $g\cdot y=y'$.

We denote the {\it orbit space} by
$$
X/G:=\{G\cdot x:x\in X\}.
$$
\end{definition}

\begin{remark}
\label{remark-not-free}
Not every finite group induces a free
action on a given space. For example,
$\mathbb{Z}/2\mathbb{Z}\approx
\{\operatorname{Id},
-\operatorname{Id}\}$
acting on all of $\mathbb{R}^2$ fixes
the origin.
\end{remark}

\section{Lifts}
\label{section-lifts}

\noindent
We begin by recalling the lifting
criterion for maps. Consider a cover
$p:E\to X$ and any map $f:Y\to X$. We
want the following diagram to commute.
$$
\xymatrix{
  & E \ar[d]^{p} \\
Y \ar@{-->}[ur]^{\tilde{f}}
  \ar[r]_{f}
  & X
}
$$

\begin{proposition}[Uniqueness of the
lift]
\label{proposition-unique-lift}
Let $p:E\to X$ be a cover and
$f:Y\to X$ a map. If two lifts
$\tilde{f}_1,\tilde{f}_2:Y\to E$ of $f$
agree at one point and $Y$ is connected,
then $\tilde{f}_1=\tilde{f}_2$.
$$
\xymatrix{
  & E \ar[d]^{p} \\
Y \ar@<0.6ex>[ur]^{\tilde{f}_1}
  \ar@<-0.6ex>[ur]_{\tilde{f}_2}
  \ar[r]_{f}
  & X
}
$$
\end{proposition}

\begin{theorem}[Homotopy lifting
property for covers]
\label{theorem-homotopy-lifting}
Given a cover $p:E\to X$, a homotopy
$f_t:Y\to X$ and a lift
$\tilde{f}_0:Y\to E$ of $f_0$, there is
a unique homotopy $\tilde{f}_t:Y\to E$
that lifts each of the $f_t$.

That is, the following diagram commutes.
$$
\xymatrix{
Y\times\{0\} \ar[r]^{\tilde{f}_0}
  \ar[dr]^{f_0}
  \ar@{^{(}->}[d]
  & E \ar[d]^{p} \\
Y\times I \ar[r]
  \ar@{-->}[ur]
  & X
}
$$
\end{theorem}

\begin{definition}
\label{definition-locally-path-connected}
We say that a topological space $Y$ is
{\it locally path-connected} if for
every point $y\in Y$ and every
neighbourhood $U$ of $y$ there is a
path-connected $V$ with
$y\in V\subseteq U$.
\end{definition}

\begin{proposition}
\label{proposition-lifting-criterion}
Let $p:(E,E_0)\to(X,x_0)$ be a cover and
$f:(Y,y_0)\to(X,x_0)$ a map, with $Y$ a
path-connected and locally
path-connected space. Then there is a
lift $\tilde{f}:(Y,y_0)\to(E,E_0)$ if
and only if
$$
f_*\pi_1(Y,y_0)\subseteq
p_*\pi_1(E,E_0).
$$
\end{proposition}

\begin{definition}
\label{definition-lift}
A {\it lift} of a map $f:Y\to X$ with
respect to a cover $p:E\to X$ is a map
$\tilde{f}:Y\to E$ such that
$p\tilde{f}=f$. That is, the following
diagram commutes.
$$
\xymatrix{
  & E \ar[d]^{p} \\
Y \ar@{-->}[ur]^{\tilde{f}}
  \ar[r]_{f}
  & X
}
$$
\end{definition}

\section{Deck transformations}
\label{section-deck-transformations}

\begin{definition}
\label{definition-deck-transformation}
Take a space $X$ and a cover
$p:E\to X$. A {\it deck transformation}
is an isomorphism of $E$ onto itself,
that is, a homeomorphism that makes the
following diagram commute.
$$
\xymatrix{
E \ar[r]^{f} \ar[d]_{p}
  & E \ar[dl]^{p} \\
X
}
$$
That is, $pf=p$.
\end{definition}

The set $\operatorname{Deck}(E)$ of all
deck transformations $f:E\to E$ has a
group structure, and we call it the
{\it group of deck transformations}, or
the {\it deck group}.

\begin{proposition}
\label{proposition-deck-agree}
If $g,f:E\to E$ are deck
transformations, with $E$
path-connected, and there is $x\in E$
with $g(x)=f(x)$, then $g=f$.
\end{proposition}

\begin{proof}
It is enough to think of a deck
transformation as a lift of the covering
projection.
$$
\xymatrix{
  & E \ar[d]^{p} \\
E \ar@{-->}[ur]^{f}
  \ar[r]_{p}
  & X
}
$$
By the lifting theorem for maps, $f$ is
unique once we choose a point
$y\in f^{-1}(x)$ in the fiber of $x$.
Note on the original: the last sentence
is unclear; the fiber of a point is its
preimage under $p$, not under $f$.
\end{proof}

\begin{definition}
\label{definition-regular-cover}
We say that a cover $p:E\to X$ is
{\it regular} if for any $x\in X$ and
$e,e'\in p^{-1}(x)$ there is a deck
transformation $f:E\to E$ such that
$f(e)=e'$. That is, the group of deck
transformations acts transitively on the
fibers.
\end{definition}

\begin{proposition}
\label{proposition-deck-normalizer}
Let $p:(E,e_0)\to(X,x_0)$ be a cover
with $E$ and $X$ path-connected and
locally path-connected. Let
$$
H=p_*\pi_1(E,e_0)\leq\pi_1(X,x_0).
$$
Then
\begin{enumerate}
\item The cover is regular if and only
if $H\trianglelefteq\pi_1(X,x_0)$.
\item There is an isomorphism
$$
\operatorname{Deck}(E)\cong N(H)/H,
$$
where
$$
N(H)=\left\{g\in\pi_1(X,x_0)
\mid g^{-1}Hg=H\right\}.
$$
\end{enumerate}
In particular, when $E$ is regular,
$$
\operatorname{Deck}(E)\cong
\pi_1(X,x_0)/H.
$$
And if moreover $E$ is the universal
cover, we simply have
$$
\operatorname{Deck}(E)\cong
\pi_1(X,x_0).
$$
\end{proposition}

\begin{proof}
The idea is to define a surjective map
$\pi_1(X,x_0)\to
\operatorname{Deck}(p)$.
Note on the original: this is all the
original says. In general such a map is
defined on $N(H)$ only.
\end{proof}

\section{Covering actions}
\label{section-covering-actions}

\begin{definition}
\label{definition-covering-action}
We say that $G\curvearrowright Y$ is a
{\it covering action} if either of the
following two equivalent conditions
holds.
\begin{itemize}
\item Every point $y\in Y$ has a
neighbourhood $U$ such that
$g_1U\cap g_2U\neq\emptyset$ implies
$g_1=g_2$.
\item Every point $y\in Y$ has a
neighbourhood $U$ such that
$gU\cap U\neq\emptyset$ implies
$g=1_G$.
\end{itemize}
In particular, these conditions imply
that the action is free.
\end{definition}

\begin{proposition}
\label{proposition-covering-action}
Given a covering action
$G\curvearrowright Y$,
\begin{enumerate}
\item The quotient projection
$p:Y\to Y/G$ given by $p(y)=Gy$ is a
regular cover.
\item If $Y$ is path-connected,
$G\approx\operatorname{Deck}(p)$.
\item If $Y$ is path-connected and
locally path-connected,
$$
G\cong\pi_1(Y/G,y_0)\big/
p_*\pi_1(Y,\tilde{y}_0).
$$
\end{enumerate}
\end{proposition}

\begin{corollary}
\label{corollary-simply-connected}
If moreover $Y$ is simply connected, by
parts (2) and (3) of
Proposition
\ref{proposition-covering-action},
$$
G=G(\tilde{Y})\cong\pi_1(Y/G,y_0).
$$
\end{corollary}

Note on the original: the notation
$G(\tilde{Y})$ is not defined. In
Section \ref{section-commented-out} the
same notation is used for a deck group.
The parts of
Proposition
\ref{proposition-covering-action} are
lettered a), b), c) in the original.

\begin{proposition}
\label{proposition-regular-quotient}
Let $p:E\to X$ be a regular cover with
$G=\operatorname{Deck}(p)$. Then there
is a homeomorphism $\psi:X\to E/G$ such
that
$$
\xymatrix{
E \ar[d]_{p} \ar[dr]^{\bar{p}} \\
X \ar[r]^{\psi} & E/G
}
$$
\end{proposition}

Note on the original: the map $\bar{p}$
is not defined and the sentence ends at
the diagram. Presumably $\bar{p}$ is the
quotient projection and the diagram
commutes.

\begin{example}
\label{example-genus-cover}
The surface of genus greater than $3$ is
a cover of $S_2$.
\end{example}

\begin{theorem}
\label{theorem-uniformization}
If $X$ is a simply connected Riemann
surface, then it is conformally
equivalent to one of the following
surfaces:
\begin{itemize}
\item $\mathbb{C}$,
\item $\mathbb{C}\cup\{\infty\}
\approx S^2$,
\item $\mathbb{H}^2$.
\end{itemize}
\end{theorem}

\begin{example}
\label{example-hyperbolic-plane}
The universal cover of the surfaces of
genus greater than or equal to $2$ is
the hyperbolic plane (recall Klein). The
cover can in fact be seen as a local
isometry, so we have a hyperbolic metric
on $S_g$.

The deck transformations are isometries,
that is,
$$
\operatorname{Deck}(p)\leq
\operatorname{Isom}^+\mathbb{H}^2
=\operatorname{PSL}(2,\mathbb{R})
=\operatorname{SL}(2,\mathbb{R})
/A\sim-A.
$$
This way we end up defining Fuchsian
groups as discrete subgroups of
$\operatorname{PSL}(2,\mathbb{R})$.
\end{example}

\begin{example}
\label{example-rotations-genus-four}
The action of $\mathbb{Z}/3\mathbb{Z}$
on $S_4$ by rotations.
\end{example}

\begin{example}
\label{example-torus-klein}
The torus with antipodal points
identified is a Klein bottle. To see it,
we take one half of the torus and look
at how the identification is on the
boundary.
\end{example}

The following three exercises are one
homework list in the original, with
items numbered $1$, $1^*$ and $2$.

\begin{exercise}[Homework]
\label{exercise-klein-lifts}
Consider the previous cover
$p:T^2\to T^2/\langle\sigma\rangle$.
Prove that every homeomorphism of the
Klein bottle onto itself admits a lift
to the torus. (In fact there are two,
one of which preserves the orientation
of the surface.)
\end{exercise}

Note on the original: $\sigma$ is not
defined; presumably it is the antipodal
map of
Example \ref{example-torus-klein}.

\begin{proof}
Building on the solution of
Exercise \ref{exercise-klein-reflection},
any homeomorphism of the Klein bottle
onto itself induces an isomorphism on
fundamental groups, and since the cover
sends $\mathbb{Z}\times\mathbb{Z}$ into
$\mathbb{Z}\rtimes\mathbb{Z}$, the
condition for lifting maps is always
satisfied.

{\bf Remark.} It seems rather that
$p_*\pi_1(T^2)$ is a subgroup of index
$2$ in $\pi_1(K)$.
\end{proof}

\begin{exercise}
\label{exercise-klein-reflection}
Prove that the homeomorphism ``reflect
with respect to $y=0$'' in the diagram
of the square of the Klein bottle lifts
to the torus.
\end{exercise}

\begin{proof}[Solution]
Let $\rho$ be that reflection. Since
the spaces are locally path-connected
and path-connected, it is enough to
show that
$$
\rho_*\pi_1(K,k_0)\subseteq
p_*\pi_1(T,t_0).
$$
The original has a figure here with
three squares whose sides carry arrows.
The first is the square of the torus:
the two vertical sides point up and the
two horizontal sides point right. Below
it, joined by an arrow $p$ pointing
down, is the square of the Klein bottle:
the two vertical sides point up, the top
side points right and the bottom side
points left. To the left of the Klein
bottle, joined to it by an arrow $\rho$,
is its reflection: the two vertical
sides point down, the top side points
left and the bottom side points right.
A dashed arrow goes from this last
square to the square of the torus.

First note that, by the theorem of Van
Kampen,
$$
\pi_1(K)\approx
\langle a,b\mid abab^{-1}\rangle
\approx\mathbb{Z}\rtimes\mathbb{Z},
$$
so we have the following diagram of
induced maps.
$$
\xymatrix{
  & \mathbb{Z}\times\mathbb{Z}
    \ar[d]^{p_*} \\
\mathbb{Z}\rtimes\mathbb{Z}
  \ar[r]_{\rho_*} \ar[ur]
  & \mathbb{Z}\rtimes\mathbb{Z}
}
$$
First let us see how $\rho_*$ acts on
the generators: $a\mapsto a^{-1}$ and
$b\mapsto b^{-1}$, so $\rho_*$ is an
isomorphism. To know how $p_*$ acts, we
see that $a\mapsto a$ and
$b\mapsto b^{-1}$. Thus
$$
p_*\pi_1(T)=
\langle a,b^{-1}\mid ab^{-1}ab\rangle
\approx\mathbb{Z}\rtimes\mathbb{Z}.
$$
We conclude that
$$
\rho_*\pi_1(K,k_0)\approx
\mathbb{Z}\rtimes\mathbb{Z}\approx
p_*\pi_1(T,t_0).
$$
It seems that the lift reverses
orientation since, at least in terms of
the generators of the fundamental group,
it sends $a\mapsto a^{-1}$.

Note on the original: the argument is
not complete. The lifting criterion
asks for an inclusion of subgroups of
$\pi_1(K)$, which does not follow from
an isomorphism of abstract groups. Also
$p_*\pi_1(T)$ is the image of the
abelian group
$\mathbb{Z}\times\mathbb{Z}$, so it is
abelian, in agreement with the remark
in the solution of
Exercise \ref{exercise-klein-lifts}.
\end{proof}

\begin{exercise}
\label{exercise-quotient-genus}
If we take $S_g$ instead of the torus,
what is $S_g/\langle\sigma\rangle$? Is
it true that
$$
\chi(S_g)=
2\chi(S_g/\langle\sigma\rangle)?
$$
Recall that $\chi(S_g)=2-2g$ and
$\chi(N_g)=2-g$.
\end{exercise}

\section{Branched covers}
\label{section-branched-covers}

\begin{definition}
\label{definition-branched-cover}
Let $E$ and $X$ be two manifolds. We say
that $p:E\to X$ is a {\it branched
cover} if there is a submanifold $B$ of
codimension $2$ in $X$ such that
$$
p|_{E\backslash p^{-1}(B)}:
E\backslash p^{-1}(B)\to X\backslash B
$$
is a covering space. We say that $B$ is
the {\it branch set}, and $p^{-1}(B)$ is
the {\it set of singular points} or
{\it branch points}.
\end{definition}

\begin{example}
\label{example-power-map}
The map $z\mapsto z^n$ acting on the
disc. The point $0$ is a branch point.
\end{example}

\begin{remark}
\label{remark-ramification-index}
Branch points locally look like
$z\mapsto z^n$, and this number is
called the {\it ramification index}.
This simply means that these are Riemann
surfaces for which there are complex
coordinate charts such that the
projection looks like $z\mapsto z^n$.
\end{remark}

\begin{theorem}
\label{theorem-properly-discontinuous}
If $G\curvearrowright S$ is properly
discontinuous, $p:S\to S/G$ is a regular
branched cover.
\end{theorem}

\begin{example}
\label{example-hyperelliptic}
The torus rotating by $\pi$ about the
$x$ axis. The quotient
$T_2/\mathbb{Z}_2$ is a sphere with four
branch points. If instead we take $S_2$,
the quotient is a sphere with $6$ branch
points. In general, $S_g$ has $2g+2$
branch points.
\end{example}

\begin{example}
\label{example-triple-torus}
Rotating a triple torus about the $x$
axis and then about the $y$ axis, we
have a cover with $4$ sheets. The branch
points, which are $6$, have index $2$.
Note that
$$
(\chi(S^2)-t)=
\sum_{x\in B}|p^{-1}(x)|,
$$
where $B$ is the set of branch points.
Note on the original: $t$ is not
defined, and the formula does not have
the form of
Theorem \ref{theorem-riemann-hurwitz}.
\end{example}

\begin{theorem}[Riemann-Hurwitz]
\label{theorem-riemann-hurwitz}
Let $p:S\to S/G$ be a branched cover
with $B\subseteq S/G$ the set of branch
points. Then
$$
\begin{aligned}
\chi(S)&=|G|\chi(S/G)
-|G|\sum_{x\in B}
\left(1-\frac{1}{n_x}\right)\\
&=|G|(\chi(S/G)-|B|)
+\sum_{x\in B}|p^{-1}(x)|,
\end{aligned}
$$
where $n_x=|G_x|$ is the size of the
stabilizer of $x$.
\end{theorem}

\begin{proof}
By lifting a cell decomposition of the
base space. The vertices may fail to
lift well, and that is why we subtract a
term.
\end{proof}

\begin{remark}
\label{remark-unbranched-euler}
For an unbranched cover $p:E\to X$ of
degree $n$,
$$
\chi(E)=n\chi(X).
$$
The following theorem generalizes this
result to branched covers.
\end{remark}

\begin{theorem}[Riemann-Hurwitz in
general]
\label{theorem-riemann-hurwitz-general}
If $p:E\to X$ is a branched cover with
$n$ sheets and $B\subseteq X$ is the set
of branch points,
$$
\chi(E)=n(\chi(X)-|B|)
+\sum_{x\in B}|p^{-1}(x)|.
$$
And if
$p^{-1}(B)=\{e_1,\ldots,e_r\}\subset E$
and the ramification index of $e_i$ is
$d_i$, then
$$
\chi(E)=n\chi(X)-\sum_{i=1}^r(d_i-1).
$$
\end{theorem}

\begin{corollary}[Hurwitz inequality]
\label{corollary-hurwitz-inequality}
If a finite group $G$ acts on $S_g$ for
$g\geq2$, then $|G|\leq84(g-1)$.
\end{corollary}

\begin{remark}
\label{remark-hurwitz-groups}
Is there a group with $|G|=84(g-1)$
that acts on $S_g$? Yes, they are called
Hurwitz groups.
\end{remark}

\begin{exercise}[Homework]
\label{exercise-three-branch-points}
Find a cover $\Sigma_g\to S^2$ with
exactly $3$ branch points. {\bf Hint}:
use the hyperelliptic involution
$\Sigma_g\to S^2$ with $2g+2$ branch
points and then another cover
$S^2\to S^2$ that ``sends two points to
three''.
\end{exercise}

\begin{proof}[Solution]
Consider first the case of the torus.
The hyperelliptic involution acts on the
torus to produce a sphere with four
branch points (an orbifold). We can then
let a rotation by $120$ degrees act on
this sphere, whose axis passes through
two of these four points and which
identifies the other two. The
composition of the two quotient maps is
a map from the torus to the sphere with
two points that have a single preimage
and one point that has two preimages.

To generalize this reasoning to the
surface of genus $g$ we simply have to
keep making rotations of the sphere to
identify points among the $2g+2$ until
only three are left. In the end, there
will be two points with a single
preimage, and the others will have
preimages of different cardinalities.

Note on the original: a rotation by
$120$ degrees cannot interchange two
points; the rotation about an axis that
interchanges two points is by $\pi$.
\end{proof}

\begin{example}
\label{example-at-least-three}
Take a branched cover $\Sigma_g\to S^2$.
If $g\geq1$, then the number of branch
points is greater than or equal to $3$.
\end{example}

\begin{exercise}[Homework]
\label{exercise-two-branch-points}
If $\Sigma_g\to S^2$ is a branched cover
with $2$ branch points, then $g=0$, that
is, $\Sigma_g$ is $S^2$.
\end{exercise}

\begin{proof}[Solution]
It seems analogous to the previous
example, simply substituting in the
Riemann-Hurwitz formula to reach a
contradiction.
\end{proof}

Now let us study something like the
converse of Riemann-Hurwitz.

Suppose that we have two orientable
surfaces $\Sigma_g$ and $\Sigma_h$. Take
a natural number $d$ (which is like the
number of sheets of the cover) and
partitions of this number (which are
like the ramification indices of each of
the $k$ branch points at its $\ell_i$
preimages):
$$
d=(d_{11},\ldots,d_{1\ell_1}),\ldots,
(d_{k1},\ldots,d_{k\ell_k}),
$$
where $\sum_id_{ji}=d$. Suppose moreover
that
$$
\chi(\Sigma_g)=d\chi(\Sigma_h)
-\sum_{i=1}^k\left(
\sum_{j=1}^{\ell_i}(d_{ij}-1)\right).
$$
Under what conditions is there a cover
$\Sigma_g\to\Sigma_h$ with $d$ sheets,
those branch points and those
ramification indices?
\begin{itemize}
\item If $h\geq1$.
\item If $h=0$, sometimes it can be
done. {\bf Homework:} what conditions on
the partitions of $d$ are needed in this
case?
\end{itemize}
Note on the original: the first item is
an unfinished sentence.

\section{Mapping class group}
\label{section-mapping-class-group}

\begin{definition}
\label{definition-finite-type}
A surface of genus $g$ with $n$
punctures and $b$ boundary components is
a {\it surface of finite type} and is
denoted by $S^b_{g,n}$.
\end{definition}

\begin{definition}
\label{definition-cross-cap}
A {\it cross-cap} is removing a disc
from a surface and gluing a M\"obius
band along the boundary of that disc.
\end{definition}

\begin{exercise}
\label{exercise-cross-caps}
What is the genus of a torus with a
cross-cap? It is a sphere with three
cross-caps, genus $3$. And a torus with
three cross-caps is a sphere with five
cross-caps, and it is also a double
torus with one cross-cap.
\end{exercise}

The point is that we can represent any
nonorientable surface $N_g$ as a sphere
with $g$ cross-caps.

\begin{definition}
\label{definition-dehn-twist}
Consider the annulus
$\mathbb{A}=S^1\times[0,1]$. The map
$$
\begin{aligned}
T:\mathbb{A}&\to\mathbb{A}\\
(e^{2\pi i\theta},t)&\mapsto
(e^{2\pi i\theta t},t)
\end{aligned}
$$
is called a {\it Dehn twist}.
\end{definition}

Note on the original: the formula is
not well defined, because $\theta$ and
$\theta+1$ give the same point but
different images. The usual Dehn twist
sends $(e^{2\pi i\theta},t)$ to
$(e^{2\pi i(\theta+t)},t)$.

\begin{remark}
\label{remark-dehn-twist-boundary}
The map $T$ is the identity on the two
ends of the annulus.
\end{remark}

\begin{definition}
\label{definition-homeo}
For a surface $\Sigma$, we denote by
$\operatorname{Homeo}(\Sigma)$ its
{\it group of homeomorphisms}.
\end{definition}

\begin{remark}
\label{remark-compact-open}
It is a topological group with the
compact-open topology. A basis (or
subbasis?) of this topology is given by
the sets of the form
$$
V(K,U)=\{f\in
\operatorname{Homeo}(\Sigma):
f(K)\subset U\}
$$
for $K$ compact and $U$ open in
$\Sigma$.
\end{remark}

\begin{definition}
\label{definition-isotopy}
A homotopy between homeomorphisms such
that all the maps in between are
homeomorphisms is called an
{\it isotopy}. If $f$ and $g$ are
isotopic we write $f\simeq g$.
\end{definition}

We denote
$$
\operatorname{Homeo}_0(\Sigma)
=\{f\in\operatorname{Homeo}(\Sigma):
f\simeq\operatorname{Id}_\Sigma\}
$$
and
$$
\operatorname{Homeo}
(\Sigma;\partial\Sigma)
=\{f\in\operatorname{Homeo}(\Sigma):
f|_{\partial\Sigma}
=\operatorname{Id}|_{\partial\Sigma}\},
$$
and we denote by
$\operatorname{Homeo}^+
(\Sigma;\partial\Sigma)$
the set of those
$f\in\operatorname{Homeo}
(\Sigma;\partial\Sigma)$
that preserve orientation.

\begin{exercise}[Homework]
\label{exercise-homeo-zero-normal}
$\operatorname{Homeo}_0(\Sigma)
\triangleleft
\operatorname{Homeo}(\Sigma)$.
\end{exercise}

\begin{proof}
Suppose that $f\simeq\operatorname{Id}$
and $g\in\operatorname{Homeo}(\Sigma)$.
Then there is a continuous map
$H:I\times\Sigma\to\Sigma$ such that
$H(0,p)=p$ and $H(1,p)=f(p)$ for
$t\in I$ and $p\in\Sigma$. Then
$$
\begin{aligned}
\tilde{H}:I\times\Sigma&\to\Sigma\\
(t,p)&\mapsto g\circ H\circ g^{-1}(p)
\end{aligned}
$$
is such that
$$
H(0,p)=g(H(0,g^{-1}(p)))
=gg^{-1}(p)=p
$$
and
$$
H(1,p)=g(f(g^{-1}(p)))
=gfg^{-1}(p),
$$
that is,
$gfg^{-1}\simeq\operatorname{Id}$.

Note on the original: the left-hand
sides of the last two displays are the
values of $\tilde{H}$, not of $H$, and
the formula for $\tilde{H}(t,p)$ is
meant as $g(H(t,g^{-1}(p)))$.
\end{proof}

\begin{definition}
\label{definition-mapping-class-group}
The {\it mapping class group} or
{\it modular group} of a surface
$\Sigma$ is
$$
\operatorname{Mod}
(\Sigma;\partial\Sigma)
=\operatorname{Homeo}
(\Sigma;\partial\Sigma)\Big/
\operatorname{Homeo}_0
(\Sigma,\partial\Sigma).
$$
\end{definition}

\begin{remark}
\label{remark-mapping-class-group}
\begin{itemize}
\item The mapping class group is a
topological group with the quotient
topology.
\item Also,
$\operatorname{Mod}(S_{g,n}^b)$ is the
quotient of
$\operatorname{Homeo}
(\Sigma;\partial\Sigma)$
by isotopy.
\item An isotopy can be thought of as a
path in the group of homeomorphisms,
that is,
$\hat{H}:I\to
\operatorname{Homeo}(\Sigma,\partial)$,
so each isotopy class is a connected
component of this group:
$$
\operatorname{Mod}(S_{g,n}^b):=
\pi_0(\operatorname{Homeo}
(\Sigma;\partial\Sigma)).
$$
\item A surface of genus $g$ with $n$
punctures and $b$ boundary components is
a {\it surface of finite type} and is
denoted by $\Sigma^b_{g,n}$. It will
turn out that if $\Sigma$ is a surface
of finite type, then
$\operatorname{Mod}
(\Sigma;\partial\Sigma)$
is a discrete group (homework).
\end{itemize}
\end{remark}

\begin{example}[Mapping class group of
the disc]
\label{example-mod-disc}
$\operatorname{Mod}
(\mathbb{D}^2,\partial\mathbb{D}^2)
\approx\{1\}$.
\end{example}

\begin{proof}
Any map of the disc that fixes the
boundary is isotopic to the identity.
Let us define the isotopy geometrically.
We look for a map
$H:\mathbb{D}^2\times I\to\mathbb{D}^2$.
The map that sends $(x,t)$ to
$$
\begin{cases}
(1-t)f(\frac{x}{1-t})
&\qquad0\leq|x|\leq1-t\\
x&\qquad1-t\leq|x|\leq1
\end{cases}
$$
is such that $H(-,0)=f$ and
$H(-,1)=\operatorname{Id}_{\mathbb{D}^2}$.

The geometric idea is, in the solid
cylinder $\mathbb{D}^2\times I$, to
collapse the disc that is one of the
ends to a point in the opposite end. In
fact, at each height we have a circle
given by $|x|=1-t$. It {\bf remains} to
see that we have a homeomorphism at
each $t$.
\end{proof}

\begin{example}[Mapping class group of
the annulus]
\label{example-mod-annulus}
$\operatorname{Mod}
(\mathbb{A};\partial\mathbb{A})
\approx\mathbb{Z}$.
\end{example}

\begin{proof}
We can think of $\mathbb{A}$ as a
surface of genus $0$ (that is, $S^2$)
with two boundary components and no
punctures.

The universal cover of $\mathbb{A}$ is
$I\times\mathbb{R}$.

Take
$[f]\in\operatorname{Mod}
(\mathbb{A},\partial\mathbb{A})$,
and $\alpha$ the curve that joins
$(0,0)$ and $(0,1)$ by a straight line.
Then $f(\alpha)$ is a curve that may
wind around the cylinder a number of
times, but ends at the same point. A
lift $\widetilde{f(\alpha)}$ is uniquely
determined by its image at a single
point. I think that this forces the lift
to be a translation of the strip
$\mathbb{R}\times I$ by as many units as
the number of times the curve
$f(\alpha)$ winds around the cylinder.
$$
\xymatrix{
  & \mathbb{R}\times I \ar[d]^{p} \\
\mathbb{A}
  \ar[ur]^{\widetilde{f(\alpha)}}
  \ar[r]_{f}
  & \mathbb{A}
}
$$
This transformation of the cover is well
defined and does not depend on the
representative. Thus we have built a
correspondence $[f]\mapsto k_f$ where
$k_f$ is a translation. Proving that
this map is surjective is not trivial,
but it is true.
\end{proof}

\begin{example}[Mapping class group of
the M\"obius band]
\label{example-mod-mobius}
$\operatorname{Mod}
(\mathbb{M};\partial\mathbb{M})
=\{\operatorname{Id}\}$.
\end{example}

\begin{example}[Mapping class group of
the sphere, homework]
\label{example-mod-sphere}
$\operatorname{Mod}(S^2)
=\{\operatorname{Id}\}$.
\end{example}

\begin{proof}
Hint: use that any homeomorphism can be
homotoped to a homeomorphism that has a
fixed point.
\end{proof}

\begin{example}[Mapping class group of
the torus]
\label{example-mod-torus}
$\operatorname{Mod}(T^2)\approx
\operatorname{SL}_2(\mathbb{Z})$.
\end{example}

\begin{proof}
Through a map
$$
\Phi:\operatorname{Mod}(T^2)\to
\operatorname{SL}_2(\mathbb{Z})
$$
that sends $[f]$ to
$$
\begin{aligned}
f_*:H_1(T^2)&\to H_1(T^2)\\
\langle a\rangle\times\langle b\rangle
&\to
\langle a\rangle\times\langle b\rangle
\end{aligned}
$$
which turns out to be surjective and
injective. At some point we use the
cover of the torus, the plane.
\end{proof}

Note on the original: the mapping class
group was defined above with all
homeomorphisms. The statements for the
sphere and the torus hold for
orientation-preserving homeomorphisms; a
reflection of $S^2$ is not isotopic to
the identity, and an
orientation-reversing homeomorphism of
$T^2$ acts on $H_1(T^2)$ with
determinant $-1$.

\begin{example}[Mapping class group of
the projective plane]
\label{example-mod-projective-plane}
$\operatorname{Mod}(\mathbb{R}P^2)
\approx\{\operatorname{Id}\}$.
\end{example}

\begin{proof}
By cutting the upper hemisphere of the
sphere along a parallel.
\end{proof}

\section{The Birman-Hilden theorem}
\label{section-birman-hilden-theorem}

\noindent
First we observe that we can take
isotopy classes and homotopy classes
interchangeably, and also that one can
take diffeomorphisms instead of
homeomorphisms.

\begin{remark}
\label{remark-homotopy-isotopy}
\begin{itemize}
\item Let $\Sigma$ be any compact
surface, and let $f$ and $g$ be
homotopic homeomorphisms of $\Sigma$.
Then $f$ and $g$ are isotopic unless
they are one of the two examples
described above (on
$\Sigma=\mathbb{D}^2$ and
$\Sigma=\mathbb{A}$). In particular, if
$f$ and $g$ are orientation-preserving
then they are isotopic.
\item Let $\Sigma$ be a compact surface.
Then every homeomorphism of $\Sigma$ is
isotopic to a diffeomorphism of
$\Sigma$.
\end{itemize}
Note on the original: this remark is in
English in the original, without a
source, and the two examples it refers
to are not in these notes.
\end{remark}

Take a branched cover $p:S\to X$ with
$S$ and $X$ orientable surfaces. Let
$B\subset X$ be the set of branch points
of $p$. We will work with homeomorphisms
that fix $B$. This can be expressed as
follows, where $k=|B|$.
$$
\operatorname{Mod}(X\backslash B)
\approx\operatorname{Mod}_{k}X
$$

Given a homeomorphism of $S$, when can
we define a homeomorphism of $X$?
$$
\xymatrix{
S \ar[r]^{\tilde{f}} \ar[d]_{p}
  & S \ar[d]^{p} \\
X \ar@{-->}[r]^{f} & X
}
$$
It is enough to ask that $f$ preserve
the fibers, that is, that
$\tilde{f}(p^{-1}(x))=p^{-1}(y)$ for
some $y\in X$. Or else,
$f(\tilde{x})=f(\tilde{x}')$ for
$\tilde{x}$ and $\tilde{x}'$ in the same
fiber of some point of $X$.

Note on the original: it is the map
$\tilde{f}$ of $S$ that preserves the
fibers, and the second condition is
presumably meant as
$p\tilde{f}(\tilde{x})
=p\tilde{f}(\tilde{x}')$.

Denote by
$\mathcal{S}\operatorname{Homeo}(S)$ the
set of those
$\tilde{f}\in\operatorname{Homeo}(S)$
that preserve the fibers of $p$, and by
$\mathcal{L}\operatorname{Homeo}(X)$ the
set of those
$f\in\operatorname{Homeo}
(X\backslash B)$
that lift under $p$.

\begin{remark}
\label{remark-natural-homomorphism}
For a branched cover $p:S\to X$, there
is a {\it natural homomorphism}
$$
\begin{aligned}
\Phi:\mathcal{S}\operatorname{Homeo}(S)
&\to\mathcal{L}\operatorname{Homeo}(X)\\
\tilde{f}&\mapsto f
\end{aligned}
$$
But the converse is not always true,
since for each liftable map there are
many possible lifts.
\end{remark}

The natural projection
$$
\begin{aligned}
\pi:\operatorname{Homeo}(S)
&\to\operatorname{Mod}(S)\\
f&\mapsto[f]
\end{aligned}
$$
allows us to define
$$
\begin{aligned}
\mathcal{S}\operatorname{Mod}(S)
:&=\pi(\mathcal{S}
\operatorname{Homeo}(S))\\
&=\{[\tilde{f}]\in
\operatorname{Mod}(S):
\exists\tilde{f}'\in
\mathcal{S}\operatorname{Homeo}(S),\
\tilde{f}\simeq\tilde{f}'\}.
\end{aligned}
$$
That is, these are the isotopy classes
in the mapping class group that contain
at least one element that preserves
fibers. The {\it symmetric mapping class
group} of $S$ is precisely
$\mathcal{S}\operatorname{Mod}(S)
<\operatorname{Mod}(S)$.

Analogously, the projection
$$
\begin{aligned}
\pi:\operatorname{Homeo}
(X\backslash B)
&\to\operatorname{Mod}_k(X)\\
f&\mapsto[f]
\end{aligned}
$$
allows us to define
$$
\begin{aligned}
\mathcal{L}\operatorname{Mod}(X)
:&=\pi(\mathcal{L}
\operatorname{Homeo}(X))\\
&=\{[f]\in\operatorname{Mod}_k(X):
\exists f'\in
\mathcal{L}\operatorname{Homeo}(X),\
f\simeq f'\},
\end{aligned}
$$
which we call the {\it liftable mapping
class group} of $X$.

In which cases can we extend a map
$\phi$ as follows? Here $\Phi$ sends
$[f]$ to $[\phi(f)]$.
$$
\xymatrix{
\mathcal{S}\operatorname{Homeo}(S)
  \ar[r]^{\phi} \ar[d]_{\pi}
  & \mathcal{L}\operatorname{Homeo}(X)
    \ar[d]^{\pi} \\
\mathcal{S}\operatorname{Mod}(S)
  \ar@{-->}[r]^{\Phi}
  & \mathcal{L}\operatorname{Mod}(X)
}
$$
First we must see that given
$[\tilde{f}]\in
\mathcal{S}\operatorname{Mod}(S)$,
if
$\tilde{g},\tilde{g}'\in
\mathcal{S}\operatorname{Homeo}(S)$
are in the class of $\tilde{f}$, then
$\Phi(\tilde{g})\simeq
\Phi(\tilde{g}')$.
For this, it is enough to see that the
covering projection preserves
homotopies.

Specifically, for $H:S\times I\to S$
with $H(-,0)=\tilde{g}$ and
$H(-,1)=\tilde{g}'$ there is a
$\bar{H}:X\times I\to X$ with
$\bar{H}(-,0)=\phi(\tilde{g})$ and
$\bar{H}(-,1)=\phi(\tilde{g}')$.

One can guarantee this if
$h_t:=H(-,t)\in
\mathcal{S}\operatorname{Homeo}(S)$
for every $t$. In fact, it is enough to
show that there is a map isotopic to the
identity that preserves fibers, since
$$
\begin{aligned}
\tilde{g}\simeq\tilde{g}'
&\implies\phi(\tilde{g})
\simeq\phi(\tilde{g}')\\
\iff\tilde{g}\circ\tilde{g}^{-1}
\simeq\operatorname{Id}_S
&\iff\phi(\tilde{g}\circ\tilde{g}^{-1}
\simeq\operatorname{Id}_X)
\iff f\simeq\operatorname{Id}_S
\end{aligned}
$$
with
$f\in\mathcal{S}\operatorname{Homeo}(S)$
through a homotopy $H:S\times I\to S$
such that
$H(-,t)\in
\mathcal{S}\operatorname{Homeo}(S)
\iff\phi(f)\simeq\operatorname{Id}_X$.

Note on the original: the last display
and the sentence after it are copied as
they stand; the chain of implications is
garbled. In the diagram, the original
labels the horizontal maps $\phi$ and
$\Phi$ in this way, while the remark
above calls the upper map $\Phi$.

\begin{definition}
\label{definition-birman-hilden-property}
Let $p:S\to X$ be a branched cover. We
say that $p$ satisfies the
{\it Birman-Hilden property} if for
every
$g\in\mathcal{S}\operatorname{Homeo}(S)
\cap\operatorname{Homeo}_0(S)$
there is an isotopy $H:S\times I\to S$
such that
$h_t:=H(-,t)\in
\mathcal{S}\operatorname{Homeo}(S)$.
That is, there is an isotopy to the
identity of $S$ that preserves the
fibers of the cover.
\end{definition}

\begin{remark}
\label{remark-well-defined}
A cover $p$ satisfies the Birman-Hilden
property if and only if the natural map
$\Phi:\mathcal{S}\operatorname{Mod}(S)
\to\mathcal{L}\operatorname{Mod}(X)$
is well defined.
\end{remark}

\begin{example}
\label{example-torus-fails}
Take $S=T^2$ as a cover of $S^2$ through
the hyperelliptic involution, which has
four branch points. Consider also
$\tilde{f}$, the rotation about the $Z$
axis by $\pi$, which is homotopic to the
identity. Moreover, $\tilde{f}$
preserves fibers; why? This cover does
not satisfy the Birman-Hilden condition.
\end{example}

\begin{theorem}[Birman-Hilden]
\label{theorem-birman-hilden}
Let $p:S\to X$ be a regular branched
cover with a finite number of sheets,
with $\chi(S)<0$. If $p$ is unbranched
or {\it solvable}
($p_*(\pi_1(S\backslash B))
\triangleleft\pi_1(X\backslash B)$
is solvable), then $p$ has the
Birman-Hilden property.
\end{theorem}

Note on the original: the condition in
parentheses is copied as it stands. The
set $B$ is a subset of $X$, not of $S$,
and it is not said which group is
required to be solvable.

We add this comment quickly. The
following definition is from
\cite[p.~464]{rotman-abstract-algebra}.
A {\it normal series} of a group $G$ is
a sequence of groups
$$
G=G_0\geq G_1\geq G_2\geq\ldots
\geq G_t=\{1\}
$$
with each $G_{i+1}$ a normal subgroup of
$G_i$; the {\it factor groups} of this
series are the quotients
$$
G_0/G_1,\quad G_1/G_2,\quad\ldots
\quad G_{t-1}/G_t.
$$
A finite group $G$ is {\it solvable} if
$G=\{1\}$ or $G$ has a normal series
whose factor groups are of finite order.

Note on the original: as written, every
finite group satisfies the last
condition. The usual definition asks
for factor groups of prime order.

Before proving it, we will prove the
following.

\begin{theorem}
\label{theorem-birman-hilden-fixed}
Let $p:S\to X$ be a regular branched
cover with a finite number of sheets $g$
and a finite number of branch points.
Suppose that each deck transformation
fixes the set of branch points and that
$S$ is hyperbolic. Then $p$ satisfies
the Birman-Hilden condition.
\end{theorem}

We begin with the following.

\medskip\noindent
{\bf Step 1.} Let $B\subset X$ be the
set of branch points. If
$f\in\mathcal{S}\operatorname{Homeo}(S)
\cap\operatorname{Homeo}_0(S)$,
then $f(\tilde{x})=\tilde{x}$ for any
$\tilde{x}\in p^{-1}(B)$.

\begin{definition}
\label{definition-riemann-surface}
A {\it Riemann surface} $S$ is a surface
for which the changes of coordinates are
analytic in the complex sense, that is,
$$
\partial_zf=0=\frac{1}{2}
(\partial_xf-i\partial_yf).
$$
\end{definition}

Note on the original: the
Cauchy-Riemann equation is the
vanishing of
$\frac{1}{2}
(\partial_xf+i\partial_yf)$,
the derivative with respect to
$\bar{z}$.

\begin{remark}
\label{remark-analytic-cover}
If $p:S\to X$ is a branched cover, we
can think of Riemann surface structures
such that $p:(S,Y)\to(X,Z)$ is analytic
in the complex sense. In this context,
the deck transformations are analytic.
\end{remark}

\begin{exercise}[Homework]
\label{exercise-trivial-center}
\begin{enumerate}
\item If $\chi(S)<0$, then
$Z(\pi_1(S))$ is trivial. (This item is
marked ``Dani'' in the original.)
\item If $f\in\operatorname{Homeo}(S)$
is such that
$f_*:\pi_1(S)\to\pi_1(S)$ is the
identity, then
$f\simeq\operatorname{Id}_S$.
\item If the map is moreover analytic,
then in fact $f$ is the identity, that
is, if
$f\in\operatorname{Homeo}(S)\cap
\operatorname{Homeo}_0(S)$
is analytic, $f=\operatorname{Id}_S$.
\end{enumerate}
\end{exercise}

\begin{proof}
This is written under the first item in
the original. {\it Another idea, more
geometric:}

First observe that a surface without
boundary $S$ with $\chi(S)<0$ admits a
Riemannian metric of constant negative
curvature; see
\url{https://mathoverflow.net/questions/14548/finding-constant-curvature-metrics-on-surfaces-without-full-power-of-uniformiza}.
To see it, it is enough to show that any
surface of genus $g$ can be seen
topologically as a polygon with $4g$
sides suitably identified.

In the case of the torus, we can use a
square and identify its sides as usual.
The key observation is that the
Euclidean plane can be tessellated with
copies of this square, and the group of
translations that sends one square to
another, $\mathbb{Z}\times\mathbb{Z}$,
is a group of isometries of the plane.

The hyperbolic plane has the incredible
property that it can be tessellated by
regular polygons with any number of
sides, in particular with a $4g$-gon
whose sides are identified to build our
surface. The group of hyperbolic
translations that takes copies of the
polygon to other copies is a subgroup of
isometries of the hyperbolic plane. In
fact, one can choose a group $\Gamma$
with that property such that the
quotient $\mathbb{H}^2/\Gamma$ is
homeomorphic to our surface.

From here on the argument is as in the
previous proof.

These ideas are also in \cite{ler}.

Note on the original: there is no
previous proof in the notes, so the
argument is incomplete.
\end{proof}

\begin{lemma}
\label{lemma-fixed-class}
Let
$f\in\operatorname{Homeo}(S)
\backslash\{\operatorname{Id}\}$
be such that $f$ fixes a point
$\tilde{x}\in S$. Consider
$$
f_*:\pi_1(S,\tilde{x})\to
\pi_1(S,\tilde{x}).
$$
Then
$$
f_*([\gamma])=[\gamma]\implies
[\gamma]=1.
$$
\end{lemma}

\begin{proof}
Let
$[\gamma]\in\pi_1(S,\tilde{x})
\backslash\{1\}$
be such that $f_*([\gamma])=[\gamma]$.
We have that
$$
f(\gamma)\simeq\gamma
\operatorname{rel}\tilde{x}.
$$
We also have the universal cover
$q:\mathbb{H}^2\to S$, and we can lift
$f$:
$$
\xymatrix{
(\mathbb{H}^2,\xi) \ar[r]^{\tilde{f}}
  \ar[d]_{q}
  & (\mathbb{H}^2,\xi) \ar[d]^{q} \\
(S,\tilde{x}) \ar[r]^{f}
  & (S,\tilde{x})
}
$$
Since $f$ is analytic,
$\tilde{f}\in
\operatorname{Isom}^+(\mathbb{H}^2)$.
Taking $\gamma:I\to S$ with
$\gamma(0)=\gamma(1)=\tilde{x}$, we
have:
$$
\xymatrix{
  & (\mathbb{H}^2,\xi) \ar[d]^{q} \\
(I,0) \ar[r]^{\gamma}
  \ar@{-->}[ur]^{\tilde{\gamma}}
  & (S,\tilde{x})
}
$$
That is, we have
$\tilde{\gamma}:I\to\mathbb{H}^2$ such
that $\tilde{\gamma}(0)=\xi$. Call
$\xi'=\tilde{\gamma}(1)\neq\xi$.

On the other hand,
$f(\gamma)\simeq\gamma
\operatorname{rel}\tilde{x}$,
so
$\widetilde{f(\gamma)}(1)
=\tilde{\gamma}(1)=\xi'$.
But
$\widetilde{f(\gamma)}
=\tilde{f}\tilde{\gamma}$
and thus
$$
\tilde{f}\xi'
$$
Note on the original: the proof stops
here, with the word ``missing''. The
proof uses that $f$ is analytic, which
is not among the hypotheses of the
lemma.
\end{proof}

\begin{lemma}
\label{lemma-commutes-with-deck}
Let
$\tilde{g}\in
\mathcal{S}\operatorname{Homeo}(S)
\cap\operatorname{Homeo}_0(S)$.
Then $\tilde{g}$ commutes with the deck
transformations.
\end{lemma}

\begin{proof}
Let $\tau\in\operatorname{Deck}(p)$, so
that
$$
\xymatrix@C=1.2em{
S \ar[rr]^{\tau} \ar[dr]_{p}
  && S \ar[dl]^{p}
     \ar[r]^{\tilde{g}^{-1}}
  & S \ar[r]^{\tau} \ar[d]^{p}
  & S \ar[dl]^{p} \ar[r]
  & S \ar[d]^{p} \\
  & X \ar[rr]_{g^{-1}}
  && X \ar[rr]^{g}
  && X
}
$$
Defining
$$
r:=\tilde{g}\circ\tau\circ
\tilde{g}^{-1}\in
\operatorname{Deck}(p)
$$
we have that:
\begin{enumerate}
\item $r\simeq\operatorname{Id}_S$.
\item $r$ fixes the branch points of the
cover $p:S\to X$.
\end{enumerate}
Given a branch point $\tilde{x}_1\in S$,
we have that
$$
r_*:\pi_1(S,\tilde{x}_1)\to
\pi_1(S,\tilde{x}_1)
$$
is an inner automorphism of
$\pi_1(S,\tilde{x}_1)$, so there is
$[\beta]\in\pi_1(S,\tilde{x}_1)$ (it
remained to see that $[\beta]\neq1$)
such that
$r_*([\gamma])
=[\beta\cdot\gamma\cdot\beta^{-1}]$.

Now if $r\neq\operatorname{Id}_S$,
since $r$ is analytic and fixes a point,
by the previous lemma,
$$
r_*([\beta])
=[\beta\cdot\beta\cdot\beta^{-1}]
=[\beta]\implies[\beta]=1.
$$
Then $r_*=\operatorname{Id}_S$, and in
fact, since this is a hyperbolic
surface, $r=\operatorname{Id}$. To
conclude, we simply note that
$$
\tilde{g}\tau\tilde{g}^{-1}\tau^{-1}
=\operatorname{Id}_S\implies
\tilde{g}\tau=\tau\tilde{g}
$$
for any $\tau\in\operatorname{Deck}(p)$.

Note on the original: the last arrow of
the upper row of the diagram has no
label. The map $r$ is defined as
$\tilde{g}\tau\tilde{g}^{-1}$, but the
conclusion uses
$\tilde{g}\tau\tilde{g}^{-1}\tau^{-1}$;
presumably $r$ is meant to be this
commutator.
\end{proof}

\begin{lemma}
\label{lemma-fixes-branch-points}
Let
$\tilde{g}\in
\mathcal{S}\operatorname{Homeo}(S)
\cap\operatorname{Homeo}_0(S)$
with a homotopy $G:S\times I\to S$ such
that $G(\tilde{x},0)=\tilde{x}$ and
$G(\tilde{x},1)=\tilde{g}(\tilde{x})$.
\begin{enumerate}
\item $\tilde{g}(\tilde{x}_i)
=\tilde{x}_i$
for $i=1,\ldots,n$, where
$p^{-1}(B)
=\{\tilde{x}_1,\ldots,\tilde{x}_n\}$.
\item The loop $\beta_i=G(\tilde{x},-)$
is homotopic to the constant loop in
$\pi_1(S,\tilde{x}_i)$.
\end{enumerate}
\end{lemma}

\begin{proof}
For the first item, by contradiction,
suppose that
$\tilde{g}(\tilde{x}_i)\neq\tilde{x}_i$,
so that, since we are assuming that the
deck transformations fix the set of
branch points,
$\tilde{g}(\tilde{x}_i)=\tilde{x}_j$
with $i\neq j$.

Take a loop $\gamma$ based at
$\tilde{x}_i$ and
$\tau\in\operatorname{Deck}(p)
\backslash\{\operatorname{Id}\}$.
By the previous lemma, $\tilde{g}$
commutes with $\tau$.
\begin{equation}
\label{equation-commute}
\tilde{g}\tau=\tau\tilde{g}
\end{equation}
Define $\beta(t)=G(\tilde{x},t)$. Then
$$
\gamma\simeq\beta\cdot\tilde{g}(\gamma)
\cdot\beta^{-1}
$$
since
$\tilde{g}\simeq\operatorname{Id}
\operatorname{rel}\tilde{x}_i$.

Applying $\tau$ to the above (since
$\tau(\tilde{x}_i)=\tilde{x}_i$),
\begin{equation}
\label{equation-apply-tau}
\tau(\gamma)\simeq\tau(\beta)\cdot
\tau\tilde{g}(\gamma)\cdot
\tau(\beta^{-1})
\end{equation}
Since $\gamma$ is arbitrary, put
$\tau(\gamma)$ and
\begin{equation}
\label{equation-replace-gamma}
\tau(\gamma)\simeq\beta\cdot
\tilde{g}(\tau(\gamma))\cdot\beta^{-1}
\end{equation}
By (\ref{equation-commute}),
$\tau(\gamma)\simeq\beta\cdot
\tau\tilde{g}(\gamma)\cdot\beta^{-1}$,
and by (\ref{equation-apply-tau}),
$$
\begin{aligned}
\tau(\beta)\cdot\tau\tilde{g}(\gamma)
\cdot\tau(\beta^{-1})
=\beta\cdot\tau\tilde{g}(\gamma)
\cdot\beta^{-1}\\
\implies(\beta^{-1}\cdot\tau(\beta))
\cdot\tau\tilde{g}(\tau)\cdot
(\tau(\beta^{-1})\cdot\beta)
\simeq\tau\tilde{g}(\gamma)
\end{aligned}
$$
Since $\gamma$ was arbitrary, we can
replace
$\gamma=t^{-1}g^{-1}(\gamma)$.

This is in order to prove that the loop
$\beta^{-1}\cdot\tau(\beta)$ commutes
with all the elements of
$\pi_1(S,\tilde{x}_j)$. Then
$$
[\beta^{-1}\cdot\tau(\beta)]\in
Z(\pi_1(S,\tilde{x}_j))\approx
\{\operatorname{Id}\}.
$$
Then
$$
\beta^{-1}\cdot\tau(\beta)\simeq
c_{\tilde{x}_j}.
$$
Hence
$$
\tau(\beta)\simeq\beta
\operatorname{rel}
\tilde{x}_i,\tilde{x}_j.
$$
{\bf Claim.} $\tau=\operatorname{Id}_S$.
$$
\xymatrix{
(\mathbb{H}^2,\xi_i)
  \ar[r]^{\tilde{\tau}} \ar[d]_{q}
  & (\mathbb{H}^2,\xi_i) \ar[d]^{q} \\
(S,\tilde{x}_i) \ar[r]_{\tau}
  & (S,\tilde{x}_i)
}
$$
Consider $\tilde{\beta}=\xi_i$, so that
$\tilde{\beta}(1)=\xi_j\in
q^{-1}(\tilde{x}_j)$.
$$
\widetilde{\tau(\beta)}
=\tilde{\tau}(\tilde{\beta}(1))
=\tilde{\tau}(\xi_j)
=\tilde{\tau}(\xi_j)=\xi_j
$$
Then $\tilde{\tau}$ fixes two points, so
$\tau=\operatorname{Id}_{\mathbb{H}^2}$,
and then $\tau=\operatorname{Id}_S$.

This claim leads us to a contradiction,
since we assumed that
$\tau\in\operatorname{Deck}
\backslash\{\operatorname{Id}\}$.

Note on the original: several formulas
of this proof are copied as they stand.
In the last aligned display
$\tau\tilde{g}(\tau)$ is presumably
$\tau\tilde{g}(\gamma)$; the
replacement
$\gamma=t^{-1}g^{-1}(\gamma)$ is
presumably by
$\tau^{-1}\tilde{g}^{-1}(\gamma)$;
$\tilde{\beta}=\xi_i$ presumably means
$\tilde{\beta}(0)=\xi_i$; and the map
that equals
$\operatorname{Id}_{\mathbb{H}^2}$ is
$\tilde{\tau}$. In the statement, the
point $\tilde{x}$ in $\beta_i$ is
presumably $\tilde{x}_i$. The second
item of the lemma is not proved. In its
place, the original has the following
two items.
\begin{enumerate}
\item Homework: prove that $\pi_1(S)$
with $\chi(S)<0$ is torsion-free.
\item Consider the diagram
$$
\xymatrix{
S \ar[d]_{p} \ar[r]^{\pi_1}
  & S/H \ar[dl]^{\pi_2} \\
S/G
}
$$
with $H\triangleleft G$. If $\pi_1$ and
$\pi_2$ (the quotient projections)
satisfy Birman-Hilden, then
$p:S\to S/G$ does too.
\end{enumerate}
\end{proof}

\begin{theorem}[Dehn-Nielsen-Baer]
\label{theorem-dehn-nielsen-baer}
$$
\operatorname{Mod}^\pm(S)\approx
\operatorname{Aut}\pi_1(S,x_0)\Big/
\operatorname{img}(\pi_1(S,x_0))
$$
\end{theorem}

\begin{proof}
$$
\xymatrix@C=1.5em{
\pi_1(S,x_0) \ar[r]^{f_*}
  & \pi_1(S,f(x_0)) \ar[r]^{c_\beta}
  & \pi_1(S,x_0) \\
[\gamma] \ar@{|->}[r]
  & [f(\gamma)] \ar@{|->}[r]
  & [\beta\circ f(\gamma)
    \circ\beta^{-1}]
}
$$
\end{proof}

Note on the original: the subgroup
written
$\operatorname{img}(\pi_1(S,x_0))$ is
presumably the group of inner
automorphisms, and the proof consists
of the diagram only.

\begin{theorem}
\label{theorem-solvable-projects}
Let $p:S\to X$ be a regular cover,
either branched or unbranched, with at
most a finite number of branch points.
Suppose that the group of deck
transformations
$G=\operatorname{Deck}(p)$ is solvable
and finite. If
$\tilde{g}\in
\mathcal{S}\operatorname{Homeo}(S)
\cap\operatorname{Homeo}_0(S)$,
then the homeomorphism $g:X\to X$ such
that
$$
\xymatrix{
S \ar[r]^{\tilde{g}} \ar[d]_{p}
  & S \ar[d]^{p} \\
X \ar[r]^{g} & X
}
$$
is isotopic to the identity on $X$, that
is, $g\in\operatorname{Homeo}_0(X)$.
\end{theorem}

The original calls this statement
``Theorem 1''.

\begin{lemma}
\label{lemma-prime-order}
If $p:S\to X$ is a cover with the
properties of
Theorem \ref{theorem-solvable-projects}
and such that $|G|$ is prime, then for
any $\tau\in\operatorname{Deck}(p)$ we
have $\tau(\tilde{x})=\tilde{x}$ for any
$\tilde{x}\in p^{-1}(B)$, where
$B\subseteq X$ is the set of branch
points.
\end{lemma}

\section{Material commented out in the
original}
\label{section-commented-out}

\noindent
The source contains two passages that
are disabled and do not appear in the
original document. They are recorded
here.

The first follows the definition of the
deck group in
Section
\ref{section-deck-transformations}.

\begin{example}
\label{example-deck-circle}
The cover $f_n:S^1\to S^1$ given by
$z\mapsto z^n$ has group of deck
transformations
$$
G(f_n:S^1\to S^1)
=\mathbb{Z}/n\mathbb{Z}.
$$
Moreover,
$f_{n*}\pi_1(S^1)=n\mathbb{Z}$. So
$$
G(f_n:S^1\to S^1)\cong
\pi_1(S^1)/f_{n*}\pi_1(S^1).
$$
And this is no coincidence.

Moreover, for the universal cover of
$S^1$, which is
$\exp:\mathbb{R}\to S^1$, it also holds,
since the fundamental group is trivial,
that is,
$$
G(\exp:\mathbb{R}\to S^1)\cong
\mathbb{Z}/\{0\}\cong\mathbb{Z}.
$$
\end{example}

The second is at the end of the source.

The {\it liftable} transformations are
those that satisfy the lifting property,
and they are denoted by
$\mathcal{L}\operatorname{Mod}(S)$.

\begin{remark}
\label{remark-descend}
For a cover $p:E\to X$ and a map
$\tilde{f}:E\to E$ that preserves
fibers, we can define a map
$f:X\to X$.
\end{remark}

The {\it symmetric} transformations are
automorphisms of the cover that preserve
fibers, and they are denoted by
$\mathcal{S}\operatorname{Mod}(S)$.

How can we relate
$\mathcal{L}\operatorname{Mod}(S)$ and
$\mathcal{S}\operatorname{Mod}(S)$? We
need an isotopy between two symmetric
transformations through symmetric
transformations.

{\bf Questions.}
\begin{itemize}
\item How must the cover be so that
there are maps
$$
\mathcal{S}\operatorname{Mod}(E)\to
\mathcal{L}\operatorname{Mod}(X),
\qquad
\mathcal{L}\operatorname{Mod}(X)\to
\mathcal{S}\operatorname{Mod}(E)?
$$
\item Under what conditions do the
equalities
$$
\mathcal{L}\operatorname{Mod}(X)
=\operatorname{Mod}(X),
\qquad
\mathcal{S}\operatorname{Mod}(E)
=\operatorname{Mod}(E)
$$
hold?
\item For the braid group $B_k$,
\end{itemize}
Note on the original: the last question
is unfinished.

\bibliography{my}
\bibliographystyle{amsalpha}

\end{document}
